\section{Surrogate models}
\label{sec:models}

Direct feedforward models map $[b,q,h_0]$ to all $50$ frames in one pass. ConvLSTM uses the same static inputs with recurrent predicted-frame feedback, initially zero. Seeded window models instead receive two static and two observed/predicted elevation channels and predict five frames per block; their conditional task is evaluated separately.

All architectures are adapted to the shared benchmark task; the original training setups are not reproduced.

\subsection{Fourier Neural Operator}
\label{subsec:fno}

FNO provides global spectral mixing at each layer \citep{li2021fno}. Figure~\ref{fig:fno-arch} shows the baseline: two coordinate channels augment the input, a $1\times1$ convolution lifts it to width $48$, and six Fourier blocks feed a two-layer pointwise head.

Each block combines spectral and local paths:
\begin{equation}
z \mapsto \mathrm{GELU}\!\left(\mathcal{K}(z) + W z\right).
\end{equation}
The spectral path applies a real 2D FFT, retains $14\times14$ modes, multiplies by learned complex weights, and transforms back. Six-cell zero padding is removed after the blocks. Mode weights can be reused on other grids, but this does not imply resolution-independent accuracy.

The baseline is trained separately on each reference. Hydrostatic FNO-modes8 and FNO-modes20 change only the retained mode count.

\begin{figure}[pos=!htbp]
\centering
\begin{tikzpicture}[
    font=\small,
    box/.style={draw, rounded corners=2pt, minimum height=8mm, minimum width=11mm, align=center, inner sep=3pt},
    lift/.style={box, fill=blue!10},
    block/.style={box, fill=orange!15},
    head/.style={box, fill=green!12},
    io/.style={box, fill=gray!10},
    op/.style={draw, circle, inner sep=1pt, minimum size=6mm, fill=orange!20},
    sm/.style={box, fill=orange!8, minimum height=7mm, minimum width=9mm},
    >={Stealth[length=2mm]},
    arr/.style={->, thick},
]

\node[io] (in) {$X$\\\scriptsize $3{\times}H{\times}W$};
\node[lift, right=7mm of in] (lift) {Lift\\\scriptsize $1{\times}1$ conv\\\scriptsize $\to d{=}48$};
\node[block, right=7mm of lift] (b1) {Fourier\\block $1$};
\node[block, right=5mm of b1] (b2) {Fourier\\block $2$};
\node[right=4mm of b2] (dots) {$\cdots$};
\node[block, right=4mm of dots] (bL) {Fourier\\block $6$};
\node[head, right=6mm of bL] (proj) {Proj.\\\scriptsize head};
\node[io, right=6mm of proj] (out) {$\widehat{Y}$\\\scriptsize $50{\times}H{\times}W$};

\draw[arr] (in) -- (lift);
\draw[arr] (lift) -- (b1);
\draw[arr] (b1) -- (b2);
\draw[arr] (b2) -- (dots);
\draw[arr] (dots) -- (bL);
\draw[arr] (bL) -- (proj);
\draw[arr] (proj) -- (out);

\node[align=center, below=2mm of lift, text width=24mm, font=\scriptsize] {+\,2 grid coord\\channels.};

\node[draw, dashed, rounded corners, fit={(b1) (bL)}, inner sep=4pt] (fitblk) {};
\coordinate (dl) at ($(fitblk.south)+(0,-3mm)$);

\node[sm] (vin) at ($(dl)+(-4.7,-1.6)$) {$z$};

\node[op, right=6mm of vin] (fft) {$\mathcal{F}$};
\node[sm, right=5mm of fft] (modes) {keep\\\scriptsize $14{\times}14$\\\scriptsize modes};
\node[op, right=5mm of modes] (rmul) {$R$};
\node[op, right=5mm of rmul] (ifft) {$\mathcal{F}^{-1}$};
\node[draw, circle, inner sep=0pt, minimum size=5mm, right=5mm of ifft, fill=red!15] (sum) {$+$};
\node[op, right=5mm of sum, fill=red!20] (act) {$\sigma$};
\node[sm, right=5mm of act] (vout) {$z'$};

\node[op, below=7mm of rmul, fill=blue!15] (W) {$W$};

\draw[arr] (vin) -- (fft);
\draw[arr] (fft) -- (modes);
\draw[arr] (modes) -- (rmul);
\draw[arr] (rmul) -- (ifft);
\draw[arr] (ifft) -- (sum);
\draw[arr] (sum) -- (act);
\draw[arr] (act) -- (vout);
\draw[arr] (vin) |- (W);
\draw[arr] (W) -| (sum);

\draw[densely dotted] (fitblk.south) -- ($(fft)!0.5!(rmul)+(0,0.95)$);


\node[align=left, font=\scriptsize, right=2mm of vout, text width=20mm]
  {$\sigma=\mathrm{GELU}$\\spectral $+$ local};
\end{tikzpicture}
\caption{Baseline FNO: three input fields plus two coordinates, width $48$, six Fourier blocks, and $50$ output frames. Each block combines a truncated $14\times14$ spectral transform with a local $1\times1$ convolution and GELU.}
\label{fig:fno-arch}
\end{figure}

\FloatBarrier
\subsection{Operator variants}
\label{subsec:operator-variants}
F-FNO uses separable one-dimensional Fourier transforms \citep{tran2023ffno}, reducing per-block spectral parameter complexity from $\mathcal{O}(d^2m_1m_2)$ to $\mathcal{O}(d^2(m_1+m_2))$.

U-FNO adds multiscale U-Net skip connections \citep{wen2022ufno}; WNO uses multiresolution wavelets \citep{tripura2023wno}. All variants use the same input features, loss, and evaluation metrics.

\subsection{Windowed rollout variants}
\label{subsec:window-models}
To assess long-horizon autoregressive drift, we train two seeded window models—Window-FNO-5 and Window-F-FNO-5—on the hydrostatic reference data. Let $Y_j$, $j = 0,\ldots,49$, denote the 50 target frames. A training window starting at frame index $s$ maps input $X_s^{\mathrm{win}} = [b, q, Y_s, Y_{s-1}]$ to target block $Y_s^{\mathrm{win}} = [Y_{s+1}, \dots, Y_{s+5}]$, setting $Y_{-1} = Y_0$ for the initial step. Window starts $s \in \{0, 5, \dots, 40, 44\}$ yield ten five-frame targets per 50-frame trajectory; the clamped final window overlaps the preceding target block by one frame to remain in range.

During inference, given the initial observation $Y_0$ as a rollout seed, the model predicts sequential five-frame blocks autoregressively by feeding its final two predictions into subsequent iterations. Evaluation metrics cover predicted frames $Y_1,\ldots,Y_{49}$ and exclude the seed. These windowed models measure conditional error accumulation and are not like-for-like direct forecasts.

\subsection{Convolutional and recurrent baselines}
\label{subsec:conv-baseline}
Convolutional and recurrent baselines test local, multiscale, and temporal representations.

The CNN uses width-$48$ $3\times3$ convolutions with GELU and a $1\times1$ readout, without pooling or global mixing.

The U-Net uses a single encoder--decoder level, average pooling, bilinear upsampling, skip connections, and base width $32$ \citep{ronneberger2015unet}.

ConvLSTM encodes $[b,q,h_0]$ with two layers to obtain static context and initial hidden/cell states. Two recurrent layers with $48$ hidden channels and $3\times3$ kernels use the usual input, forget, output, and candidate gates \citep{shi2015convlstm}. At each step, context is combined with a $1\times1$ projection of the previous predicted frame; no true target frame is supplied. A $1\times1$ readout emits each surface frame and supplies feedback for the next step, initially zero.

\subsection{Ensemble uncertainty}
\label{subsec:model-ensemble}
Seven independently seeded FNOs form the deep ensemble \citep{lakshminarayanan2017deepensembles}. Their mean gives the prediction and their per-element spread measures seed-dependent optimization variability, rather than observation noise (Section~\ref{subsec:uncertainty}).
